\documentclass[letterpaper, 11pt]{article}
\usepackage[utf8]{inputenc}
\usepackage[spanish]{babel}
\usepackage[T1]{fontenc}
\usepackage{charter}
\usepackage{csquotes}
\usepackage[margin=3cm]{geometry}
\usepackage{fancyhdr}
\usepackage{minted}
\usepackage{hyperref}
\usemintedstyle{vs}

% --- DEFINICIÓN DE ESTILOs DE PÁGINA ---
\fancypagestyle{sections}{
    \fancyhf{} 
    %\fancyhead[L]{Proyecto Final LTII} --> Texto a la izquierda
    \fancyhead[R]{Planteamiento de secciones}
    \fancyfoot[C]{\thepage}
    \setlength{\headheight}{15pt}
    \renewcommand{\headrulewidth}{0.4pt}
}

\fancypagestyle{technic}{
    \fancyhf{} 
    \fancyhead[R]{Detalle técnico del desarrollo}
    \fancyfoot[C]{\thepage}
    \setlength{\headheight}{15pt}
    \renewcommand{\headrulewidth}{0.4pt}
}
% -----------------------------


\title{Proyecto Final Laboratorio Temático II}
\author{@Char[50] \textit{Carlos T. Mendoza}}
\date{\today}

\renewcommand{\thesection}{\Alph{section}}
\renewcommand{\thesubsection}{\roman{subsection}}

\begin{document}
\maketitle
\tableofcontents

\newpage
\section{Antecedentes}
    Steam ha existido desde \textit{septiembre de 2003} y aunque en sus inicios tuvo un arranque algo \enquote{atropellado}, hoy en día es la plataforma más enorme y poderosa en el ámbito del \textit{Gaming} para \textit{PC}. Incluso ahora \enquote{y gracias al desarollo Proton \footnote{Capa intermedia que traduce a Linux las instrucciones de un videojuego nativo de Windows.}} ha incursionado en el mundo de los \enquote{PC's de mano} con su \textbf{SteamDeck} y, a últimas fechas, en el de \enquote{consolas} con la mítica \textbf{SteamMachine}. 

    Más allá de lo asombroso que pueda ser, es inegable que ostenta un monopolio extremadamente sobrado \textit{\enquote{controlando al rededor de un 70 a 75\% de la distribución de videojuegos en formato digital para PC}}. Por lo anterior, mi intención con este proyecto no es sólo apoyarme de sus datos para ilustrar algunos tópicos interesantes, sino también poner en contexto al espectador tanto de lo \enquote{bueno} como lo \enquote{malo} de la plataforma y su influencia. 

    Los datos utilizados para este proyecto fueron seleccionados tanto por su fuente \enquote{puesto que vienen directamente de la plataforma en cuestion} como por su relevancia en el ámbito comercial y del \textit{Gaming}. 


\section{Contenido}
    \subsection{Sección A} %Texto introductorio y mencion a Gabe Newell% 
    Antes de entrar en materia y hablar a profundidad de Steam, es imposible dejar de lado a su empresa matriz y, aún más importante, a su fundador: \textbf{Gabe Newell}. Nacido el 3 de noviembre de 1962 en Colorado, E.U. Siempre fué un hombre de visión excepcional y entusiasta de los videojuegos. 

\vspace{0.5cm}

    Tuvo la oportunidad de estudiar en Harvard pero su estadía no duraría lo suficiente como para terminar la carrera, puesto que \textit{aproximadamente en 1983} empezaría a trabajar en Microsoft trás la sugerencia de \textbf{Steve Ballmer} \enquote{CEO en aquel tiempo} de \enquote{hacer algo útil}. Dentro de dicha compañia, Newell pasaría los próximos 13 años de su vida desarrollando y siendo parte de las primeras versiones de Windows. Hasta que en 1996, junto con un compañero de trabajo, \textbf{Mike Harrington}, tomaría la decisión de fundar Valve Corporation con la intención de seguir su pasion por los videojuegos. 

\vspace{0.5cm}

    Durante sus primeros pasos, Valve desarollaría algunos de los videojuegos más legendarios de la época, tales como \enquote{Half-lyfe} o el tan afamado \enquote{Counter-Strike}. Sin embargo, la forma en la que se hacian las actualizaciones \footnote{Los usuarios debían descargar e instalar cada parche por su cuenta.} era un verdadero dolor de cabeza y necesitaban una alternativa. Por tal motivo, a mediados de 2002, Valve les presentaría un proyecto a empresas gigantezcas como Microsoft, Yahoo! y Amazon para desarrollar una plataforma en conjunto.   

\vspace{0.5cm}

    Desgraciadamente, ninguna de estas empresas mostro el suficiente interés o compromiso con el proyecto. 


\newpage
    \subsection{Sección B} %El origen de Steam% 
    Aproximadamente un año más tarde Valve lanzaria oficialmente \textbf{Steam} y, apenas días después, \textbf{Counter-Strike 2}. Pero el lanzamiento no sería de forma aislada, de hecho se volvió obligatorio instalar Steam para poder realizar la a Counter-Strike 2, lo que ayudo a la compañia a crear una base sólida de usuarios desde su salida. 

    De esta manera lograban dos objetivos, brindarle a sus jugadores un cliente que se encargara de las tediosas actualizaciones y, al mismo tiempo, planteaban un ambiente único y propio; cuestión que más adelante le daría un impulso mucho más alto a la plataforma. Pues lo que comensaría como un simple cliente, pronto se convertiría en un espacio donde los usuarios podían interactuar entre ellos, llevar el registro de sus logros e incluso poder disfrutar de otros juegos además de los desarrollados por Valve.

    \subsection{Sección C} %VALVE en el mercado actual% 
    
% -1- Su monopolio: 
% Datos de mercado sobre el control y distribución digital según la plataforma. 
% (Puede ser desde que empezo a distribuir (2005 o 2006 aprox.) y el 2025 para ver el contraste || ver la evolución año con año hasta hoy).

    \subsection{Sección D}
% -2- Indie vs AA vs AA: 
% Datos sobre el "ÉXITO" de cada tier por año, basado en {Score Methacritic y Ventas Totales}. 
% (Al ser nociones diferentes habría que NORMALIZARLO)* {recuerda que tienes una nota de voz sobre este punto} 

    \subsection{Sección E}
% -3- Transición de preferencia de Hardware y Software (Windows -> Linux):
% Datos cuantitativos sobre el uso de hardware y software en función de ilustrar el cambio de Windows a Linux en conjunto con el desarrolo de Proton.
% (Dependiendo de la representación puede mostrarse de una u otra manera la supremacía de cada sistema y hardware) 

    \subsection{Sección F}
% -3- El pozo sin fondo de las bibliotecas: 
% Datos sobre la posesión exagerada de títulos por usuario. 
% (Podría ser un análisis por contraste o evolución, como en el primer punto)


\newpage
\pagestyle{technic}
\section{Aspectos técnicos}
\subsection{Objetivo 1}
    Puesto que los datos recabados de cada plataforma no incluyen un valor explícto de las ventas e ingresos para cada una, fue necesarío realizar estimaciones que me permitíeran realizar las gráficas. 


\end{document}






